\section{Execution and feedback verification}
\label{sec:checks}
Tool execution and memory updates require verifiable records of what actually ran.
The execution controller computes cryptographic hashes over input data, executed code,
and returned allocations. If an agent modifies an artifact after execution, previous
validation receipts become invalid. The controller returns \texttt{PASS} if all required
guards pass, \texttt{FAIL} if a guard detects a defect, and \texttt{INCOMPLETE} if
evidence is missing. This mechanism enforces execution provenance without modifying the
underlying algorithm.

For cross-episode memory, timing and accounting checks verify that delayed outcomes
cleared settlement before they update persistent state. Trades that rely on look-ahead
data or unadjusted corporate splits are blocked, preventing corrupted feedback from
distorting the agent's memory. Unsettled trades wait for settlement, while valid
unprofitable trades remain eligible learning signals.

The audit study evaluates models under unconstrained generation, optional self-checking,
and mandatory verification. Appendix~\ref{app:audit} describes the evaluation setup,
defect detection benchmarks, and empirical audit results across model tiers.
